\documentclass[11pt]{article}
\usepackage[T1]{fontenc}
\usepackage[utf8]{inputenc}
\usepackage{lmodern}
\usepackage[margin=1in]{geometry}
\usepackage{amsmath,amssymb,amsthm}
\usepackage{parskip}
\usepackage{microtype}
\usepackage{booktabs,array}
\usepackage{hyperref}
\hypersetup{colorlinks=true,linkcolor=blue,citecolor=blue,urlcolor=blue,
  pdftitle={Congruences for an Iterative Generating Function: A396798},
  pdfauthor={Research note prepared for Kaoru Aguilera Katayama}}
\newtheorem{theorem}{Theorem}
\newtheorem{lemma}[theorem]{Lemma}
\newtheorem{corollary}[theorem]{Corollary}
\theoremstyle{remark}
\newtheorem{remark}[theorem]{Remark}
\newcommand{\Z}{\mathbb Z}
\newcommand{\eps}{\varepsilon}
\newcommand{\itn}[2]{{#1}^{\circ #2}}
\newcommand{\coeff}[2]{[x^{#1}]\,#2}
\title{Congruences for an Iterative Generating Function:\\
  A Proof of the A396798 Iteration Conjectures}
\author{Kaoru Aguilera Katayama}
\date{September 22, 2026}

\begin{document}
\maketitle

\begin{abstract}
Let $A(x)\in x+x^2\Z[[x]]$ be determined by
$A(x)=x+\itn{A}{4}(x)\itn{A}{5}(x)$, where superscripts with a
composition symbol denote iteration. We prove that, for every integer $k$,
\[
\itn{A}{k}(x)\equiv
\frac{x}{1-kx}+4\eps_k\frac{x^4}{(1-x)^3}\pmod 8,
\qquad \eps_k\in\{0,1\},\quad \eps_k\equiv k\pmod 2.
\]
This establishes the seven iteration congruences listed as conjectures in
OEIS A396798 at the time of consultation. It also gives the precise starting
index of the coefficient pattern for $A$ itself. The proof uses uniqueness
of a formal fixed point and an explicit rational function over
$\Z/8\Z$. In this ring the reduced series has compositional order exactly
eight. An accompanying dependency-free Python program checks the rational
identities exactly and independently reconstructs initial coefficients.
\end{abstract}

\section{The problem and its documented status}

Paul D. Hanna's sequence A396798 is defined in OEIS by the ordinary
generating function
\begin{equation}\label{eq:definition}
A(x)=\sum_{n\geq1}a(n)x^n
     =x+\itn{A}{4}(x)\itn{A}{5}(x).
\end{equation}
Thus $a(1),a(2),\ldots$ begin $1,1,9,133,2517,55609$.
The notation $\itn{A}{j}$ means $j$ successive substitutions of $A$,
not the ordinary product power $A(x)^j$. We set $\itn{A}{0}(x)=x$.
Negative iterates use the compositional inverse, which exists because
the coefficient of $x$ is one.

The source record \cite{oeis}, checked on September 22, 2026, labels
seven assertions about iterates $2$ through $8$ as conjectures. They
predict vanishing coefficient tails for even iterates and periodic
residue blocks for odd iterates. The result below proves all seven.
The record also proposes a residue block for $a(n)$ itself; its stated
starting index needs the correction explained in Section~\ref{sec:consequences}.
This is a specific set of published sequence conjectures. The source
check documents their listed status; it does not certify that no
unindexed proof exists elsewhere.

Throughout, congruences of formal series mean coefficientwise congruences.
No analytic convergence is required. Put $R=\Z/8\Z$ and define
\begin{equation}\label{eq:model}
M_k(x)=\frac{x}{1-kx},\qquad
H(x)=\frac{x^4}{(1-x)^3},\qquad
F(x)=M_1(x)+4H(x).
\end{equation}
Every denominator here has constant coefficient one and is therefore
invertible in the relevant formal power series ring.

\begin{theorem}\label{thm:main}
The defining equation~\eqref{eq:definition} has a unique solution in
$x+x^2\Z[[x]]$. For this solution and every $k\in\Z$,
\begin{equation}\label{eq:main}
\itn{A}{k}(x)\equiv M_k(x)+4\eps_kH(x)\pmod 8,
\end{equation}
where $\eps_k$ is the parity of $k$. In particular, for $n\geq2$,
\begin{equation}\label{eq:coefficients}
\coeff{n}{\itn{A}{k}(x)}
\equiv k^{n-1}+4\eps_k\binom{n-2}{2}\pmod 8,
\end{equation}
where $\binom{0}{2}=\binom{1}{2}=0$. The coefficient of $x$ is one
for every $k$.
\end{theorem}

\section{Uniqueness of the defining formal series}

The following argument works over a commutative ring, including rings
with zero divisors. This matters when reducing modulo eight.

\begin{lemma}\label{lem:fixedpoint}
Let $S$ be a commutative ring with identity and let $r,s\geq1$ be
integers. There is exactly one $B\in x+x^2S[[x]]$ satisfying
$B=x+\itn{B}{r}\itn{B}{s}$.
\end{lemma}

\begin{proof}
For $B\in x+x^2S[[x]]$ write
$T(B)=x+\itn{B}{r}\itn{B}{s}$. Suppose $B-C\in x^NS[[x]]$, with
$N\geq2$. Substitution by a series of order one preserves divisibility
by $x^N$. Also, if $U-V$ is divisible by $x^N$, then so is
$D(U)-D(V)$ for every formal series $D$ for which the substitutions
are defined: apply $U^j-V^j=(U-V)\sum_{i=0}^{j-1}U^{j-1-i}V^i$
term by term. Induction on the number of substitutions consequently gives
$\itn{B}{j}-\itn{C}{j}\in x^NS[[x]]$ for all $j\geq0$.

Both factors in the product defining $T$ have order one, so
\begin{align*}
T(B)-T(C)
&=(\itn{B}{r}-\itn{C}{r})\itn{B}{s}
  +\itn{C}{r}(\itn{B}{s}-\itn{C}{s})\\
&\in x^{N+1}S[[x]].
\end{align*}
Starting with $B_0=x$ and putting $B_{j+1}=T(B_j)$, we obtain
$B_{j+1}-B_j\in x^{j+2}S[[x]]$. Hence each coefficient eventually
stabilizes. Their coefficientwise limit is a fixed point because addition,
multiplication, and these substitutions respect finite truncations.

If two fixed points first differed in degree $N$, their images under
$T$ could not differ before degree $N+1$, a contradiction. This proves
uniqueness. The construction also shows that reducing coefficients
modulo an integer commutes with taking this fixed point.
\end{proof}

\section{An explicit series of order eight}

The elementary fractional-linear identity
\begin{equation}\label{eq:mobius}
M_u\bigl(M_v(x)\bigr)=M_{u+v}(x)
\end{equation}
holds over any commutative ring. The additional term $4H$ in $F$
has a useful property: products of two such perturbations vanish in $R$.

\begin{lemma}\label{lem:square}
In $R[[x]]$, $\itn{F}{2}(x)=M_2(x)$.
\end{lemma}

\begin{proof}
Since $4^2=0$ in $R$, the binomial theorem gives the exact first-order
substitution rule
\begin{equation}\label{eq:taylor}
P(U+4V)=P(U)+4V P'(U)
\end{equation}
for formal series $U,V\in xR[[x]]$. There are no divisions by
factorials in this identity. Apply it to $F=M_1+4H$:
\[
\itn{F}{2}=M_1\circ M_1+
4\bigl(H\circ M_1+(M_1'\circ M_1)H\bigr).
\]
It is enough to evaluate the parenthesis modulo two. In $\mathbb F_2[[x]]$,
direct substitution gives
\[
H\bigl(M_1(x)\bigr)=\frac{x^4}{1-x},
\qquad
M_1'\bigl(M_1(x)\bigr)H(x)=\frac{x^4}{1-x}.
\]
Their sum is zero. The correction multiplied by four therefore vanishes
modulo eight, and~\eqref{eq:mobius} proves the claim.
\end{proof}

\begin{lemma}\label{lem:iterates}
For every integer $k$, the identity
$\itn{F}{k}=M_k+4\eps_kH$ holds in $R[[x]]$.
\end{lemma}

\begin{proof}
First take $k\geq0$. By Lemma~\ref{lem:square} and~\eqref{eq:mobius},
$\itn{F}{2j}=M_{2j}$. For the next iterate,
\[
\itn{F}{2j+1}=F\circ M_{2j}
=M_{2j+1}+4H\circ M_{2j}.
\]
But $M_{2j}(x)\equiv x\pmod2$, so
$4H\circ M_{2j}=4H$ in $R[[x]]$. This proves the formula for
nonnegative $k$.

In particular, $\itn{F}{8}=M_8=x$. Thus $F$ has an inverse under
composition and its integer iterates depend only on $k$ modulo eight.
The right-hand side $M_k+4\eps_kH$ has the same dependence. The
formula consequently holds for negative $k$ as well.
\end{proof}

\section{Verification of the fixed point}

We now identify $F$ with the reduction of $A$. This step verifies the
defining equation itself, rather than merely matching initial coefficients.

\begin{lemma}\label{lem:equation}
The series $F$ satisfies $F=x+\itn{F}{4}\itn{F}{5}$ in $R[[x]]$.
\end{lemma}

\begin{proof}
Lemma~\ref{lem:iterates} gives
$\itn{F}{4}=M_4=x+4x^2$ and $\itn{F}{5}=M_5+4H$ in $R[[x]]$.
Consequently,
\[
\itn{F}{4}\itn{F}{5}
=\frac{x^2(1+4x)}{1-5x}+4xH.
\]
Using $F-x=x^2/(1-x)+4H$, their difference is
\begin{align*}
F-x-\itn{F}{4}\itn{F}{5}
&=\frac{x^2}{1-x}-\frac{x^2(1+4x)}{1-5x}+4(1-x)H\\
&=4x^4\left(\frac{1}{(1-x)(1-5x)}+
                    \frac{1}{(1-x)^2}\right)
\quad\text{in }R[[x]].
\end{align*}
The expression in parentheses is zero modulo two, because
$1-5x\equiv1-x\pmod2$. After multiplication by four it is zero
modulo eight. This proves the identity.
\end{proof}

\begin{proof}[Proof of Theorem~\ref{thm:main}]
Apply Lemma~\ref{lem:fixedpoint} over $\Z$ to obtain $A$, and over
$R$ to its reduction. Lemma~\ref{lem:equation} shows that $F$ is
the latter fixed point. Uniqueness therefore gives $A\equiv F\pmod8$.
Composition, including composition with inverses, respects coefficient
reduction. Lemma~\ref{lem:iterates} proves~\eqref{eq:main}.

Finally,
$M_k(x)=\sum_{n\geq1}k^{n-1}x^n$ and
$H(x)=\sum_{n\geq4}\binom{n-2}{2}x^n$.
Extracting coefficients gives~\eqref{eq:coefficients} and the separate
coefficient-one statement in degree one. For $k=0$, one may simply use
$M_0=x$, avoiding any need to interpret $0^0$.
\end{proof}

\section{The conjectures and the indexing correction}\label{sec:consequences}

For odd $k$, the coefficient tail beginning at $n=2$ repeats the block
$(k,1,k+4,5)$ modulo eight. Indeed, $k^2\equiv1\pmod8$, while the
parity of $\binom{n-2}{2}$ repeats $0,0,1,1$ from $n=2$ onward.
For even $k$, the perturbation vanishes and $k^3\equiv0\pmod8$.
The complete results for one period of iterates are as follows.

\begin{center}
\renewcommand{\arraystretch}{1.15}
\begin{tabular}{@{}cl@{}}
\toprule
Iterate & Exact description modulo $8$\\
\midrule
$1$ & Coefficient of $x$: $1$; from $n=2$: $(1,1,5,5)$ repeating\\
$2$ & $x+2x^2+4x^3$\\
$3$ & Coefficient of $x$: $1$; from $n=2$: $(3,1,7,5)$ repeating\\
$4$ & $x+4x^2$\\
$5$ & Coefficient of $x$: $1$; from $n=2$: $(5,1,1,5)$ repeating\\
$6$ & $x+6x^2+4x^3$\\
$7$ & Coefficient of $x$: $1$; from $n=2$: $(7,1,3,5)$ repeating\\
$8$ & $x$\\
\bottomrule
\end{tabular}
\end{center}

Rows $2$ through $8$ prove the seven iteration assertions in the source.
For the remaining coefficient claim, the repeating block $(1,1,5,5)$
starts at $n=2$, with $a(1)=1$ separate. Starting that same block at
$n=1$ would predict $a(3)\equiv5$, whereas $a(3)=9\equiv1$.
Thus the literal starting index in the consulted record cannot be retained.
The corrected formula, valid for every index, is
\[
a(1)=1,\qquad
a(n)\equiv
\begin{cases}
1 & n\equiv2,3\pmod4,\\
5 & n\equiv0,1\pmod4
\end{cases}
\quad(n\geq2).
\]

\begin{corollary}
The reduction of $A$ modulo eight has compositional order exactly eight.
For all $k\in\Z$, $\itn{A}{k+8}\equiv\itn{A}{k}\pmod8$.
\end{corollary}

\begin{proof}
The eighth iterate is $x$. Conversely, the coefficient of $x^2$ in
$\itn{A}{k}$ is $k$ modulo eight, so an iterate equal to $x$ must
have $8\mid k$.
\end{proof}

\section{Reproducible exact checks}

The accompanying file \texttt{verify\_a396798.py} uses Python's standard
library only. It verifies rational functions by cross-multiplying their
polynomial numerators and denominators in $(\Z/8\Z)[x]$. All denominators
have odd constant coefficients, so these checks certify identities of
formal series to every degree. In particular, the script checks
$F\circ F=M_2$, each transition between the eight proposed iterate
formulas, and the defining identity $F=x+\itn{F}{4}\itn{F}{5}$.

As separate arithmetic checks, it reconstructs the first twenty integer
coefficients directly from~\eqref{eq:definition} and compares them with
the source data. It then reconstructs $A$ modulo eight through degree
$64$ and checks every coefficient of iterates $k=0,\ldots,16$ against
the theorem. All checks pass. These finite comparisons supplement the
formal proof; the rational identities and the uniqueness argument supply
the statements for arbitrary degree.

Run the checks with \texttt{python3 verify\_a396798.py} in GITHUB \url{https://github.com/PolloXDDD/a396798_repository}. The manuscript
is self-contained and compiles with \texttt{pdflatex}; no external
figures or bibliography files are required.


\begin{thebibliography}{9}
\bibitem{oeis}
Paul D. Hanna,
\emph{A396798: G.f. satisfies $A(x)=x+A^4(x)A^5(x)$},
The On-Line Encyclopedia of Integer Sequences, June 16, 2026.
Here the superscripts in the source title denote iteration.
\url{https://oeis.org/A396798}; internal record:
\url{https://oeis.org/A396798/internal}.
Consulted September 22, 2026.
\end{thebibliography}

\end{document}
